\chapter{Transactions and Isolation}
\label{ch:transactions}

A \term{transaction} is a group of reads and writes that the database treats as
one unit. Transactions let an application reason about one operation at a time,
even though many run concurrently and any of them may be interrupted by a crash.

\section{ACID}

\begin{description}[leftmargin=1.2em]
\item[Atomicity.] All of a transaction's writes happen, or none do. A transaction that aborts, or is cut off by a crash, leaves no trace. (Provided by the log: \cref{ch:wal,ch:aries}.)
\item[Consistency.] A transaction takes the database from one valid state to another, where ``valid'' is defined by the application's invariants. The engine helps with constraints, but consistency is mostly the application's promise.
\item[Isolation.] Concurrent transactions do not see each other's intermediate states; the result is as if they ran in some order. (Provided by concurrency control: \cref{ch:concurrency}.)
\item[Durability.] Once a commit is acknowledged, its effects survive any later crash. (Provided by the log and \code{fsync}.)
\end{description}

\section{Schedules and serializability}

A \term{schedule} is the interleaved order in which the operations of
concurrent transactions actually execute. A schedule is \term{serializable} if
its effect equals that of \emph{some} serial order of the same transactions.
Serializability is the gold standard of isolation: the application cannot tell
that concurrency happened.

Two operations \term{conflict} if they come from different transactions, touch
the same item, and at least one is a write. A schedule is
\term{conflict-serializable} if it can be turned into a serial one by swapping
adjacent non-conflicting operations. The test is a graph: draw an edge
$T_i \to T_j$ for every pair of conflicting operations where $T_i$'s comes
first. The schedule is conflict-serializable exactly when the graph has no cycle.

\section{Anomalies}

Weaker isolation admits specific bad interleavings. Naming them is the clearest
way to define the isolation levels \cite{berenson}.

\begin{description}[leftmargin=1.2em]
\item[Dirty read.] $T_2$ reads a value $T_1$ wrote but has not committed. If $T_1$ aborts, $T_2$ read data that never existed.
\item[Non-repeatable read.] $T_1$ reads an item, $T_2$ changes it and commits, and $T_1$ reads it again and sees a different value.
\item[Phantom.] $T_1$ reads the set of items matching a condition (``every student''), $T_2$ inserts a new matching item and commits, and $T_1$ repeats the query and sees an extra row.
\item[Lost update.] Both transactions read a counter, both add one, both write: one increment disappears.
\end{description}

\section{Isolation levels}

\begin{table}[ht]
\centering
\begin{tabular}{@{}lccc@{}}
\toprule
Level & Dirty read & Non-repeatable read & Phantom \\
\midrule
Read uncommitted & possible & possible & possible \\
Read committed & --- & possible & possible \\
Repeatable read & --- & --- & possible \\
Serializable & --- & --- & --- \\
\bottomrule
\end{tabular}
\caption{The SQL standard's levels, defined by the anomalies they forbid.}
\label{tab:isolation}
\end{table}

Berenson et al.\ \cite{berenson} showed that defining levels by these phenomena
alone is ambiguous, and that real systems' levels (notably snapshot isolation,
which forbids all three but admits a ``write skew'') do not fit the table
cleanly. The table remains the right first mental model.

\begin{keyidea}
Isolation is a dial, not a switch. Lower levels allow more concurrency and
specific anomalies; serializable allows none. Each level is defined by what it
forbids, and each forbidden anomaly corresponds to a specific lock held for a
specific time (\cref{ch:concurrency}).
\end{keyidea}

\begin{ontodb}[: transactions]
Every shell statement is its own transaction unless \code{\textbackslash begin}
opens an explicit one, ended by \code{\textbackslash commit} or
\code{\textbackslash abort}. Isolation is a setting,
\code{\textbackslash set isolation ru|rc|rr|ser}, defaulting to read committed.
The isolation contract is written as executable schedules: twelve specs, each a
scripted interleaving of two sessions with expectations such as ``session 2 is
blocked'', ``the read returns Dana'', or ``session 2 is aborted''. They cover
dirty reads, non-repeatable reads, lost updates, lock upgrades, two deadlock
shapes, write--write conflicts, locks released on abort, and phantoms.

In a persistent database today, each update statement (or \code{\textbackslash
load}) is one atomic, durable transaction (\cref{ch:durability}).
\end{ontodb}

\begin{status}
Statement-level atomicity and durability run today. Explicit multi-statement
transactions and isolation levels are plans 7.2 to 7.6.
\end{status}

\section*{Further reading}
Gray and Reuter \cite{grayreuter}; Berenson et al.\ \cite{berenson}; Petrov
\cite{petrov}, chapter 5.
